\section{Related Work}
\label{sec:related_work}
The confluence of generative models and engineering design has been a burgeoning area of research. In particular, the use of deep learning for the generation of 3D shapes has seen significant advancements. Early approaches, such as Generative Adversarial Networks (GANs) and Variational Autoencoders (VAEs), have been applied to 3D object generation, but often struggle with producing high-resolution and structurally complex shapes.

More recently, diffusion models \cite{ho2020denoising, nichol2021improved} have emerged as a powerful class of generative models, capable of producing high-fidelity images and 3D shapes. The application of diffusion models to 3D shape generation is a rapidly evolving field.

The concept of hierarchical representation has also been explored in the context of 3D shape generation. By representing shapes in a hierarchical manner, it is possible to capture both coarse and fine-grained details, leading to more realistic and detailed geometries. Our work builds upon these ideas, using a hierarchical representation to decode the latent vectors from our diffusion model.

Furthermore, the integration of physics-based simulations into the generative process is a key area of research. By incorporating feedback from simulations, it is possible to generate designs that are not only aesthetically pleasing but also physically plausible and optimized for specific performance criteria. In the context of aerospace engineering, the use of CFD simulations to guide the design of aerodynamic shapes is a well-established practice. However, the integration of CFD solvers directly into the training loop of a generative model is a challenging endeavor, due to the computational cost of the simulations. Our work addresses this challenge by leveraging a GPU-accelerated CFD solver, which enables us to efficiently evaluate the aerodynamic performance of the generated designs during training.

The Transformer architecture, first introduced in \cite{vaswani2017attention}, has revolutionized the field of natural language processing and is now being applied to a wide range of computer vision and generative modeling tasks. The attention mechanism at the core of the Transformer allows the model to capture long-range dependencies and complex relationships in the data. In our work, we draw inspiration from the Transformer architecture in the design of our latent diffusion model.

Furthermore, this work is conceptually guided by two key principles: "Thinking in Diffusion and Speaking in Autoregression" (TiDAR) and "Hierarchical Representation Mapping" (HRM). The LiDAR principle suggests a process where a vast design space is explored in parallel through a diffusion process, followed by a more structured, sequential evaluation. This mirrors the creative process of rapid ideation followed by critical analysis. The HRM principle, inspired by recursive models that demonstrate strong general reasoning on grid-based benchmarks, advocates for hierarchical and voxelized representations to effectively capture complex spatial relationships, which is particularly well-suited for structural design tasks.

